% !TeX root = ../bd-screen.tex

\ifdefined\ollangid
\chapter{导论}

模态逻辑是经典逻辑被算子 $\Box$（「盒子」）与 $\Diamond$（「钻石」）扩张的结果，这些算子作用于!!{formula}。直观上，$\Box$ 可读作「必然」，$\Diamond$ 可读作「可能」，因此 $\Box p$ 是「$p$ 必然真」，而 $\Diamond p$ 是「$p$ 可能真」。由于必然性与可能性的基本形而上学概念，模态逻辑显然具有重大的哲学意义。它允许对诸如「$\Box p \lif p$」（若 $p$ 必然，则其为真）或「$\Diamond p \lif \Box\Diamond p$」（若 $p$ 可能，则其必然可能）之类的形而上学原理加以形式化。

算子 $\Box$ 与 $\Diamond$ 是\emph{内涵的}。这意味着 $\Box !A$ 或 $\Diamond !A$ 是否成立并非只取决于 $!A$ 是否成立。不是内涵的算子则是\emph{外延的}。否定是外延的：$\lnot !A$ 当且仅当 $!A$ 不成立；所以 $\lnot !A$ 是否成立只取决于 $!A$ 是否成立。$\Box$ 与 $\Diamond$ 并非如此：$\Box !A$ 或 $\Diamond !A$ 是否成立还取决于~$!A$ 的意义。虽然通常的真值函项语义足以处理外延算子，但像 $\Box$ 与 $\Diamond$ 这样的内涵算子需要另一种语义。本书中占据中心地位的这样一种语义就是关系语义（亦称可能世界语义或\zhFirst{Kripke}{克里普克}语义）。

对于对应于把 $\Box$ 解释为「必然」的逻辑而言，这种语义相对简单：一个解释~$\mModel{M}$ 不再给!!{propositional variable}指派真值，而是给它们指派一个「世界」集合——直观上，即那些使 $p$ 被解释为真的世界~$w$。基于这样的解释，我们可以定义一个满足关系。这个满足关系的定义使得 $\Box !A$ 在世界~$w$ 处得到满足当且仅当 $!A$ 在\emph{所有}世界处得到满足：$\mSat{M}{\Box !A}[w]$ 当且仅当对所有世界~$v$ 有 $\mSat{M}{!A}[v]$。这对应于\zhFirst{Leibniz}{莱布尼茨}的观念：必然真的东西就是在每个可能世界中都真的东西。

「必然」并不是解释 $\Box$ 算子的唯一方式，但它是标准的解释——「必然」与「可能」就是所谓的\emph{真性}模态。其他解释把 $\Box$ 读作「（某人~$A$）知道」，「某人 $A$ 相信」，「应当如此」，或「将总是真的」。这些分别是认知、信念、道义与时态模态。对 $\Box$ 的不同解释会使不同的!!{formula}成为逻辑真，并认定不同的推理为有效的。例如，一切必然的东西与一切被知道的东西都是真的，所以 $\Box !A \lif !A$ 在真性解释与认知解释下都是逻辑真。相比之下，并非一切被相信的、也并非一切应当如此的都会真的如此，所以 $\Box !A \lif !A$ 在信念或道义解释下不是逻辑真。

为了处理模态算子的不同解释，语义由世界之间的一种关系——所谓的可及关系——加以扩展。于是 $\mSat{M}{\Box !A}[w]$ 当且仅当对所有从~$w$ 可及的世界~$v$ 有 $\mSat{M}{!A}[v]$。由此得到的语义非常灵活而强大，其基本思想可用来为基于其他内涵算子的逻辑提供语义解释。这样一种逻辑是直觉主义逻辑，一种基于\zhFullName{L. E. J. Brouwer}{L. E. J. 布劳威尔}{布劳威尔}的构造性数学分支的构造性逻辑。直觉主义逻辑之所以在哲学上有趣，正是因为这个原因——它在数学的构造性说明中扮演着重要角色——不过在 20 世纪它也被有影响力的英国哲学家\zhFullName{Michael Dummett}{迈克尔·达米特}{达米特}提出为优于经典逻辑的逻辑。关系模型的另一个应用是作为虚拟条件句（即反事实条件句）的语义，这一进路由\zhFullName{Robert Stalnaker}{罗伯特·斯达内克}{斯达内克}与\zhFullName{David K. Lewis}{大卫·刘易斯}{刘易斯}开创。

本书是内涵逻辑的句法、语义与证明论导论。它只处理命题逻辑，虽然未来的版本也将处理谓词逻辑。全书分为三个部分：第一部分处理正规模态逻辑。这些逻辑具有算子 $\Box$ 与~$\Diamond$。我们讨论它们的句法、关系模型以及基于它们的语义概念（如有效性与后承），以及!!{derivation}系统（既包括公理系统也包括表列）。我们建立这些逻辑的一些基本结果，如所考察的!!{derivation}系统的可靠性与完全性，并讨论诸如过滤之类的一些模型论构造。第二部分处理直觉主义逻辑。这里我们讨论自然演绎与公理化!!{derivation}、关系语义与拓扑语义，以及这些!!{derivation}系统的可靠性与完全性。第三部分处理反事实条件句的\zhFirst{Lewis-Stalnaker}{刘易斯–斯达内克}语义。附录讨论集合论与关系理论中一些对关系语义至关重要的思想与结果，并回顾经典命题逻辑的句法、语义与!!{derivation}理论。
\else
\chapter{Introduction}

Modal logics are extensions of classical logic by the operators $\Box$
(``box'') and $\Diamond$ (``diamond''), which attach to !!{formula}s.
Intuitively, $\Box$ may be read as ``necessarily'' and $\Diamond$ as
``possibly,'' so $\Box p$ is ``$p$ is necessarily true'' and $\Diamond
p$ is ``$p$ is possibly true.'' As necessity and possibility are
fundamental metaphysical notions, modal logic is obviously of great
philosophical interest. It allows the formalization of metaphysical
principles such as ``$\Box p \lif p$'' (if $p$ is necessary, it is
true) or ``$\Diamond p \lif \Box\Diamond p$'' (if $p$ is possible,
it is necessarily possible).

The operators $\Box$ and $\Diamond$ are \emph{intensional}. This means
that whether $\Box !A$ or $\Diamond !A$ holds does not just depend on
whether $!A$ holds or doesn't.  An operator which is not intensional
is \emph{extensional}. Negation is extensional: $\lnot !A$ holds iff
$!A$ does not; so whether $\lnot !A$ holds only depends on whether
$!A$ holds or doesn't. $\Box$ and $\Diamond$ are not like that:
whether $\Box !A$ or $\Diamond !A$ holds depends also on the meaning
of~$!A$.  While ordinary truth-functional semantics is enough to deal
with extensional operators, intensional operators like $\Box$ and
$\Diamond$ require a different kind of semantics. One such semantics
which takes center stage in this book is relational semantics (also
called possible-worlds semantics or Kripke semantics). 

For the logic which corresponds to the interpretation of $\Box$ as
``necessarily,'' this semantics is relatively simple: instead of
assigning truth values to !!{propositional variable}s, an
interpretation~$\mModel{M}$ assigns a set of ``worlds'' to
them---intuitively, those worlds~$w$ at which $p$ is interpreted as true.
On the basis of such an interpretation, we can define a satisfaction
relation. The definition of this satisfaction relation makes $\Box !A$ satisfied at a world~$w$ iff $!A$ is
satisfied at \emph{all} worlds: $\mSat{M}{\Box !A}[w]$ iff
$\mSat{M}{!A}[v]$ for all worlds~$v$. This corresponds to Leibniz's
idea that what's necessarily true is what's true in every possible world.

``Necessarily'' is not the only way to interpret the $\Box$ operator,
but it is the standard one---``necessarily'' and ``possibly'' are the
so-called \emph{alethic} modalities. Other interpretations read $\Box$
as ``it is known (by some person~$A$) that,'' as ``some person $A$
believes that,'' ``it ought to be the case that,'' or ``it will always
be true that.'' These are epistemic, doxastic, deontic, and temporal
modalities, respectively. Different interpretations of $\Box$ will
make different !!{formula}s logically true, and pronounce different
inferences as valid. For instance, everything necessary and everything
known is true, so $\Box !A \lif !A$ is a logical truth on the alethic
and epistemic interpretations. By contrast, not everything believed
nor everything that ought to be the case actually is the case, so
$\Box !A \lif !A$ is not a logical truth on the doxastic or deontic
interpretations. 

In order to deal with different interpretations of the modal
operators, the semantics is extended by a relation between worlds, the
so-called accessibility relation.  Then $\mSat{M}{\Box !A}[w]$ iff
$\mSat{M}{!A}[v]$ for all worlds~$v$ which are accessible from~$w$.
The resulting semantics is very versatile and powerful, and the basic
idea can be used to provide semantic interpretations for logics based
on other intensional operators. One such logic is intuitionistic
logic, a constructive logic based on L. E. J. Brouwer's branch of
constructive mathematics. Intuitionistic logic is philosophically
interesting for this reason---it plays an important role in
constructive accounts of mathematics---but was also proposed as a logic
superior to classical logic by the influential English philosopher
Michael Dummett in the 20th century. Another application of relational
models is as a semantics for subjunctive, or counterfactual,
conditionals, an approach pioneered by Robert Stalnaker and David K.
Lewis.

This book is an introduction to the syntax, semantics, and proof
theory of intensional logics. It only deals with propositional logics,
although future editions will also treat predicate logics.  The
material is divided into three parts: The first part deals with normal
modal logics. These are logics with the operators $\Box$
and~$\Diamond$. We discuss their syntax, relational models and
semantic notions based on them (such as validity and consequence) and
!!{derivation} systems (both axiomatic systems and tableaux). We establish some
basic results about these logics, such as the soundness and
completeness of the !!{derivation} systems considered, and discuss some
model-theoretic constructions such as filtrations. The second part
deals with intuitionistic logic. Here we discuss natural deduction and
axiomatic !!{derivation}s, relational and topological semantics, and
soundness and completeness of the !!{derivation} systems. The third part deals
with the Lewis-Stalnaker semantics of counterfactual conditionals. The
appendices discuss some ideas and results from set theory and the
theory of relations that's crucial to the relational semantics, and
review syntax, semantics, and !!{derivation} theory of classical propositional
logic.
\fi
